% Ad Atticum 5.11 --- headnote
% ad-atticum-05-11, 6 July 51 BC, written at Athens

\textit{Cicero to Atticus, written at Athens on 6 July 51 BC
(the day before the Nones of July), Cicero's tenth and last
day in the city before setting sail for Asia. He has been
writing letters home and none to Atticus, and opens with a
spasm of self-reproach and a plea --- ``by all our
fortunes!'' --- that Atticus prevent any prolongation of the
Cilician command. The journey-letters' refrain is now a
drumbeat: \textit{flagrem desiderio urbis}, ``I burn with
longing for the city.'' A swift political bulletin follows:
the consul M. Claudius Marcellus has had a magistrate of
Novum Comum flogged --- an act repudiating Caesar's grant of
Latin rights to Transpadane Gaul, foul to Cicero and a
deliberate provocation to Caesar --- and Pompey, with whom
Cicero has been corresponding through the Greek freedman
Theophanes, has been talked out of any flight to Spain.}

\textit{The body of the letter is the proconsul's own progress
report. He is at Athens with his legate Pomptinus, the
quaestor, and a flotilla of light Rhodian and Mytilenean
craft; only Atticus's friend Tullius is missing from the
staff. He believes the journey's good Greek reputation is
holding, with the proverbial caveat \textit{[Greek: hoiaper
hē despoina, toiai kai therapainides]} --- as the mistress,
so are the maids --- that staff discipline rests on his own.
The final two sections are pure Atticus-business: prefectures
to be handed out (Cicero will not repeat the error he made
in the case of Appuleius), the Epicurean garden's request
that he ask the Athenian Areopagus to vacate a decree against
their school, routed delicately through C. Memmius (the
patron of Lucretius, then at Mytilene, who held the lease on
the disputed property), and finally a quiet word about Pilia,
Atticus's wife: her letter to Cicero arrived very tender, do
not let her know he has read it. The corrupt
\textit{\textdagger{}nomanaria\textdagger{}} at the close has
defied emendation; the body translation preserves the cross,
and the textual decision belongs to a future translator's
note.}
